% 2.5 — L'analyse statique des commandes shell
\section{L'analyse statique des commandes shell}
\label{sec:statique}

Lire une commande avant de l'exécuter relève de l'analyse statique. Trois familles d'outils
existent.

\begin{description}[style=nextline,leftmargin=0.6cm,font=\sffamily\bfseries]
  \item[Les motifs (expressions régulières)]
    C'est l'approche des \textit{hooks} de sécurité les plus courants, et celle de la bêta
    de \demo{} : repérer \cmd{rm -rf}, \cmd{DROP TABLE}, \cmd{--force}. Elle est rapide et
    lisible, mais elle lit des caractères là où le shell lit une grammaire : citations,
    commentaires, substitutions, enveloppes, variables (tableau~\ref{tab:regex-pieges}).
  \item[Les analyseurs de qualité]
    ShellCheck signale les erreurs de syntaxe, les pièges sémantiques et les cas limites
    d'un script~\cite{shellcheck}. Son but est la \emph{justesse} du script, pas l'effet
    d'une ligne sur un système de fichiers donné.
  \item[Les analyseurs syntaxiques]
    \cmd{bashlex} est une transposition en Python de l'analyseur de GNU bash, qui produit un
    arbre syntaxique complet sans rien exécuter~\cite{bashlex} ; \cmd{tree-sitter-bash} est
    une grammaire bash pour l'outil d'analyse incrémentale tree-sitter~\cite{tree-sitter-bash}.
\end{description}

\paragraph{Ce qu'un arbre syntaxique ne dit pas.} Un arbre dit que \cmd{mv a.txt b.txt}
est une commande \cmd{mv} à deux opérandes. Il ne dit pas que \cmd{b.txt} existe, donc que
son contenu va être perdu ; ni que \cmd{-n} aurait empêché la perte ; ni que, après un
\cmd{cd build}, \cmd{rm -rf *} vise d'autres fichiers ; ni que \cmd{\$X} vaut \cmd{-rf}.
Pour passer de la \emph{structure} à l'\emph{effet}, il faut deux connaissances que
l'analyse syntaxique n'apporte pas :
\begin{enumerate}
  \item la \textbf{sémantique de chaque outil} --- sa grammaire d'options et ce que chaque
    option fait (\cmd{cp} écrase, \cmd{cp -n} n'écrase pas, \cmd{sed -i} réécrit en place,
    \cmd{gzip} supprime l'original) ;
  \item l'\textbf{état du système de fichiers}, mis à jour au fil de la ligne (ce qui
    existe, ce que la commande précédente vient de créer ou de déplacer).
\end{enumerate}
C'est ce qui conduit \demo{} à un \emph{moteur} : un analyseur, une table de règles qui
décrit la sémantique des outils, et une simulation sur une vue virtuelle du système de
fichiers (chapitre~\ref{ch:conception}).
